\section{Обобщённая структура компиляционного синтезатора}
\label{sec:structure}

\subsection{Задание 3.4. Схема синтезатора}

На рисунке~\ref{fig:structure} изображена обобщённая структура
синтезатора речи, использующего компиляционный метод, промоделированный в
данной работе. Жирными рамками выделены блоки, реализованные в нашей
<<системе>> (вырезание фонем/слогов, склейка и просодическая модификация);
пунктирными --- блоки, отсутствующие в ней.

\begin{figure}[H]
	\centering
	\resizebox{\textwidth}{!}{%
	\begin{tikzpicture}[
			node distance=0.5cm and 0.35cm,
			box/.style={rectangle, draw, rounded corners=2pt, align=center,
				font=\footnotesize, inner sep=4pt, text width=3.1cm, minimum height=1.0cm},
			boxd/.style={box, dashed},
			base/.style={rectangle, draw, rounded corners=2pt, align=center,
				font=\footnotesize, inner sep=4pt, text width=3.2cm, minimum height=0.8cm,
				fill=gray!10},
			arr/.style={-{Stealth}, thick},
		]
		\node[boxd] (text) {Орфографический текст};
		\node[boxd, right=of text] (norm) {Нормализация текста};
		\node[boxd, right=of norm] (tr) {Транскриптор \\ (правила транскрибирования)};
		\node[boxd, right=of tr] (pros) {Просодическая разметка \\ (синтагмы, ударения)};
		\node[box, below=of pros] (select) {Выбор единиц \\ (фонемы / слоги / дифоны)};
		\node[box, left=of select] (concat) {Конкатенация \\ (склейка, сглаживание стыков)};
		\node[box, left=of concat] (mod) {Просодическая модификация \\ ($F_0(t)$, $A(t)$, $T(t)$)};
		\node[box, left=of mod] (out) {Речевой сигнал};
		\node[base, below=of select, yshift=-0.2cm] (base) {База единиц речи \\ (фонетическая база)};

		\draw[arr] (text) -- (norm);
		\draw[arr] (norm) -- (tr);
		\draw[arr] (tr) -- (pros);
		\draw[arr] (pros) -- (select);
		\draw[arr] (select) -- (concat);
		\draw[arr] (concat) -- (mod);
		\draw[arr] (mod) -- (out);
		\draw[arr] (base) -- (select);
	\end{tikzpicture}%
	}
	\caption{Обобщённая структура компиляционного синтезатора речи}
	\label{fig:structure}
\end{figure}

\subsection{Каких компонентов недостаёт}

Чтобы разработанный <<синтезатор>> стал полноценной системой
автоматического синтеза речи с неограниченным словарём, ему не хватает
следующих компонентов:
\begin{enumerate}
	\item модуля анализа текста: расшифровка чисел, дат, аббревиатур,
	омографов, сокращений, работа с незнакомыми словами;
	\item автоматического транскриптора, реализующего все правила
	транскрибирования, включая редукцию безударных гласных, оглушение и
	озвончение согласных, учёт исключений;
	\item модуля просодической разметки, который по синтаксису и пунктуации
	автоматически расставляет синтагмы, словесные и фразовые ударения и
	выбирает тип интонации;
	\item развитой базы единиц речи: аллофонов и дифонов для всех возможных
	звукосочетаний, а не только для ограниченного набора фонем;
	\item алгоритма выбора единиц по критерию качества стыковки;
	\item автоматического сглаживания стыков (например, модификация по
	PSOLA) и согласования тона соседних единиц;
	\item просодической модификации по прогнозируемым, а не заданным
	вручную контурам $F_0(t)$, $A(t)$, $T(t)$.
\end{enumerate}
Кроме того, для естественного звучания важно учитывать коартикуляцию
(взаимное влияние соседних звуков) и непрерывность мелодического контура,
что достигается переходом от фонем к дифонам и применением процедур
сглаживания на стыках.
